% SGA 3，报告 VII_B，英文重建本——第 1 轮（中文译本）。
% 范围：依据原文件第 16 页，包含 0.5 全部；止于原文件第 17 页的第 1 节之前。
% 合并阅读版页码： 529. 再版印刷页码： 518.
% 原版页边定位： 491.
% 依据： Polo--Gille Exp7B-13oct24.pdf,
% SHA-256 D303690E0195654BBE7C3DAA5D89EA23E24CD709A3575217234AC19AE8836694.
% 曾参阅 Jacob Reinhold 的 CC BY 4.0 Markdown，作为对比版本来源。

\section*{0.5. 基变换}
\phantomsection
\label{sga3:VIIB:0.5}
\addcontentsline{toc}{subsection}{\SGARefLink{sga3:VIIB:0.5}{0.5}. 基变换}

设
\[
  \phi:k\longrightarrow\ell
\]
是伪紧环的同态（\SGARefLink{sga3:VIIB:0.1.3}{0.1.3}）。
\SGARefLink{sga3:VIIB:0.3}{0.3} 的构造可以作如下推广。\footnote{编者展开了本段的细节。}

\medskip
\noindent\textbf{定义 0.5.A。}
\phantomsection
\label{sga3:VIIB:0.5.A}
{\itshape 对于每个 \(M\in\operatorname{PC}(k)\) 和
\(N\in\operatorname{PC}(\ell)\)，令
\[
  M\widehat{\otimes}_kN
\]
为 \(M\otimes_kN\) 对于由下列投影的核所定义的线性拓扑的分离完备化：
\[
  M\otimes_kN
  \longrightarrow
  (M/M')\otimes_k(N/N'),
\]
其中 \(M'\) 和 \(N'\) 分别遍历 \(M\) 和 \(N\) 的开子模。
则 \(M\widehat{\otimes}_kN\) 是伪紧 \(\ell\)-模。若 \(m\in M\) 且 \(x\in N\)，
则将 \(m\otimes_kx\) 在 \(M\widehat{\otimes}_kN\) 中的典范像记为
\[
  m\widehat{\otimes}_kx.
\]}

特别地，这适用于 \(N=\ell\)。这时
\[
  M\widehat{\otimes}_k\ell
\]
称为通过基变换 \(k\to\ell\) 从 \(M\) 得到的伪紧 \(\ell\)-模。

\medskip
\noindent\textbf{注记 0.5.B。}
\phantomsection
\label{sga3:VIIB:0.5.B}

\noindent\textit{\SGARefLink{sga3:vii:paragraph:5-1:equation:i}{(i)}}
在这样的基变换中，\(\ell\) 不必是预有限 \(k\)-代数。
典型的例子是 \(k\) 为域、\(\ell=K\) 为 \(k\) 的任意域扩张的情形。

\smallskip
\noindent\textit{\SGARefLink{sga3:vii:paragraph:5-1:equation:ii}{(ii)}}
然而，若 \(N\) 的底 \(k\)-模是伪紧的，例如 \(\ell\) 是预有限 \(k\)-代数，
则 \SGARefLink{sga3:VIIB:0.4}{0.4} 表明 \(N\) 的每个开 \(k\)-子模都包含一个开
\(\ell\)-子模。这时 \(M\widehat{\otimes}_kN\) 与
\SGARefLink{sga3:VIIB:0.3}{0.3} 中伪紧 \(k\)-模 \(M\) 和 \(N\) 的完备化张量积一致，
所以记号没有歧义。

在任何情形下，通过限制标量得到的 \(k\)-模 \(N|_k\) 都是拓扑 \(k\)-模：映射
\[
  k\times N\longrightarrow N,
  \qquad
  (t,n)\longmapsto\phi(t)n
\]
是连续的。我们以 \(\operatorname{Hom}_c(M,N|_k)\) 表示从 \(M\) 到
\(N|_k\) 的连续 \(k\)-线性映射组成的阿贝尔群。

\medskip
\noindent\textbf{命题 0.5.C。}
\phantomsection
\label{sga3:VIIB:0.5.C}
\textit{对于每个 \(M\in\operatorname{PC}(k)\) 和
\(N\in\operatorname{PC}(\ell)\)，存在函子性同构}
\[
  \operatorname{Hom}_{\operatorname{PC}(\ell)}
    (M\widehat{\otimes}_k\ell,N)
  \simeq
  \operatorname{Hom}_c(M,N|_k).
\]
\footnote{因此，当 \(\ell\) 是预有限 \(k\)-代数时，函子
\(M\mapsto M\widehat{\otimes}_k\ell\) 是限制标量函子的左伴随。}

事实上，设 \(u:M\to N|_k\) 是连续同态。映射
\[
  u':M\times\ell\longrightarrow N,
  \qquad
  (m,\lambda)\longmapsto\lambda u(m)
\]
连续且“双线性”，即对第一个因子 \(k\)-线性，对第二个因子 \(\ell\)-线性。
若 \(N'\) 是 \(N\) 的开 \(\ell\)-子模，\SGARefLink{sga3:VIIB:0.3.1}{引理 0.3.1}
的论证给出开子模 \(M'\subset M\) 和开理想 \(\ell'\subset\ell\)，使得
\[
  u'(M'\times\ell)\subset N',
  \qquad
  u'(M\times\ell')\subset N'.
\]
因此，\(u\) 诱导连续的 \(\ell\)-线性映射
\[
  U:M\widehat{\otimes}_k\ell\longrightarrow N
\]
满足
\[
  U(m\widehat{\otimes}_k\lambda)=\lambda u(m).
\]
反之，每个态射
\[
  f:M\widehat{\otimes}_k\ell\longrightarrow N
\]
给出态射
\[
  f':M\longrightarrow N|_k,
  \qquad
  m\longmapsto f(m\widehat{\otimes}_k1).
\]

如同 \SGARefLink{sga3:VIIB:0.3.2}{0.3.2}、\SGARefLink{sga3:VIIB:0.3.5}{0.3.5} 和 \SGARefLink{sga3:VIIB:0.3.1.2}{0.3.1.2}，由此得到：

\medskip
\noindent\textbf{推论 0.5.D。}
\phantomsection
\label{sga3:VIIB:0.5.D}
{\itshape 函子
\[
  \operatorname{PC}(k)\longrightarrow\operatorname{PC}(\ell),
  \qquad
  M\longmapsto M\widehat{\otimes}_k\ell
\]
右正合，并与滤过投射极限交换。因此，若 \((M_i)\) 是
\(\operatorname{PC}(k)\) 中的滤过投射系，则有典范同构
\[
  (\varprojlim_iM_i)\widehat{\otimes}_k\ell
  \simeq
  \varprojlim_i(M_i\widehat{\otimes}_k\ell).
\]
此外，若 \(M\) 和 \(N\) 是伪紧 \(k\)-模，则有典范同构
\[
  (M\widehat{\otimes}_kN)\widehat{\otimes}_k\ell
  \simeq
  (M\widehat{\otimes}_k\ell)
    \widehat{\otimes}_{\ell}
  (N\widehat{\otimes}_k\ell).
\]}

\medskip
\noindent\textbf{定义 0.5.E。}
\phantomsection
\label{sga3:VIIB:0.5.E}
{\itshape 最后，若 \(A\) 是预有限 \(k\)-代数，则 \(A\widehat{\otimes}_k\ell\)
上存在唯一的预有限 \(\ell\)-代数结构，使得
\[
  (a\widehat{\otimes}_k\lambda)
  (b\widehat{\otimes}_k\mu)
  =
  ab\widehat{\otimes}_k(\lambda\mu)
\]
对 \(a,b\in A\) 及 \(\lambda,\mu\in\ell\) 成立。
称预有限 \(\ell\)-代数 \(A\widehat{\otimes}_k\ell\) 是从 \(A\) 沿
\(k\to\ell\) 通过标量扩张或基变换得到的。}

% 下一断点： Exp7B 原文件页码 17, 合并阅读版页码 530,
% 再版印刷页码 519, 第 1 节处。
